\section{Conclusion}
\label{sec:conclusion}

A frequency-aware diffusion objective can do exactly what it was designed to
do---measurably correcting a real spectral deficiency, through the component
its design credits---and still lose to the baseline it was meant to improve
on. The correction is worth something, and we measure what: $\boostValue$~FID.
The objective delivers it and gives back three times as much elsewhere, so that
a post-process obtaining the same correction for free outperforms it.

The broader lesson is about attribution. It is not enough to show that a method
changes the quantity it targets; the change must be priced against what the
evaluation metric can register, and against what the method costs to obtain it.
Both of the measurement hazards we report were caught only because we asked
what a result \emph{should} have been before accepting what it was.

\todo{revisit once the learned-weighting arms report}
